\textbf{Phần 1. Trắc nghiệm lựa chọn}
\setcounterpageref{ex}{0}

\begin{ex} %[5D10N56]
$1$ thế kỉ bằng bao nhiêu năm?
\choice
{$10$ năm}
{\True $100$ năm}
{$1,000$ năm}
{$50$ năm}
\loigiai{
$1$ thế kỉ $= 100$ năm.\\
\textbf{Đáp án của câu hỏi là B}
}
\end{ex}

\begin{ex} %[5D10N56]
Năm $2024$ thuộc thế kỉ nào?
\choice
{\True Thế kỉ XXI}
{Thế kỉ XX}
{Thế kỉ XIX}
{Thế kỉ XXII}
\loigiai{
Từ năm $2001$ đến $2100$ là thế kỉ XXI. Do đó năm $2024$ thuộc thế kỉ XXI.\\
\textbf{Đáp án của câu hỏi là A}
}
\end{ex}

\begin{ex} %[5D10N56]
Tháng nào sau đây có $31$ ngày?
\choice
{Tháng 4}
{Tháng 6}
{Tháng 9}
{\True Tháng 7}
\loigiai{
Các tháng có $31$ ngày là: Tháng 1, 3, 5, 7, 8, 10, 12.\\
\textbf{Đáp án của câu hỏi là D}
}
\end{ex}

\begin{ex} %[5D10H56]
$2{,}5$ giờ bằng bao nhiêu phút?
\choice
{\True $150$ phút}
{$120$ phút}
{$125$ phút}
{$250$ phút}
\loigiai{
$2{,}5 \times 60 = 150$ phút.\\
\textbf{Đáp án của câu hỏi là A}
}
\end{ex}

\begin{ex} %[5D10H56]
Đổi $3$ ngày $12$ giờ ra giờ ta được
\choice
{$36$ giờ}
{$72$ giờ}
{$96$ giờ}
{\True $84$ giờ}
\loigiai{
$3 \times 24 + 12 = 72 + 12 = 84$ giờ.\\
\textbf{Đáp án của câu hỏi là D}
}
\end{ex}

\begin{ex} %[5D10H56]
Đổi $216$ phút ra giờ ta được
\choice
{$3{,}36$ giờ}
{\True $3{,}6$ giờ}
{$3{,}5$ giờ}
{$3{,}4$ giờ}
\loigiai{
$216 : 60 = 3{,}6$ giờ.\\
\textbf{Đáp án của câu hỏi là B}
}
\end{ex}

\begin{ex} %[5D10H56]
$\dfrac{3}{4}$ giờ bằng bao nhiêu phút?
\choice
{\True $45$ phút}
{$30$ phút}
{$40$ phút}
{$50$ phút}
\loigiai{
$\dfrac{3}{4} \times 60 = 45$ phút.\\
\textbf{Đáp án của câu hỏi là A}
}
\end{ex}

\begin{ex} %[5D10V56]
Một năm nhuận có bao nhiêu ngày?
\choice
{$364$ ngày}
{$365$ ngày}
{\True $366$ ngày}
{$367$ ngày}
\loigiai{
Năm nhuận tháng 2 có 29 ngày nên cả năm có $366$ ngày.\\
\textbf{Đáp án của câu hỏi là C}
}
\end{ex}

\begin{ex} %[5D10V56]
Bác Hồ đọc Tuyên ngôn Độc lập khai sinh ra nước Việt Nam Dân chủ Cộng hòa vào năm $1945$. Năm đó thuộc thế kỉ nào?
\choice
{Thế kỉ XVIII}
{\True Thế kỉ XX}
{Thế kỉ XIX}
{Thế kỉ XXI}
\loigiai{
Năm $1945$ thuộc thế kỉ XX (từ năm 1901 đến năm 2000).\\
\textbf{Đáp án của câu hỏi là B}
}
\end{ex}

\begin{ex} %[5D10V56]
Đổi $2$ giờ $45$ phút ra số thập phân có đơn vị là giờ ta được
\choice
{$2{,}45$ giờ}
{$2{,}3$ giờ}
{\True $2{,}75$ giờ}
{$2{,}8$ giờ}
\loigiai{
$2$ giờ $45$ phút $= 2 + \dfrac{45}{60} = 2{,}75$ giờ.\\
\textbf{Đáp án của câu hỏi là C}
}
\end{ex}

\begin{ex} %[5D10V56]
Một trận đấu bóng đá diễn ra trong $90$ phút và nghỉ giữa hai hiệp $15$ phút. Tổng thời gian của trận đấu kể cả thời gian nghỉ là
\choice
{$1{,}5$ giờ}
{$1{,}65$ giờ}
{$2$ giờ}
{\True $1{,}75$ giờ}
\loigiai{
Tổng thời gian: $90 + 15 = 105$ phút $= 105 : 60 = 1{,}75$ giờ.\\
\textbf{Đáp án của câu hỏi là D}
}
\end{ex}

\begin{ex} %[5D10C56]
Một người làm việc từ $7$ giờ $45$ phút đến $11$ giờ $30$ phút, nghỉ trưa $1$ giờ $15$ phút rồi làm tiếp đến $16$ giờ $45$ phút. Tổng thời gian người đó thực sự làm việc trong ngày là
\choice
{$7$ giờ}
{$7$ giờ $45$ phút}
{\True $7$ giờ $30$ phút}
{$8$ giờ}
\loigiai{
Sáng làm: $11$ giờ $30$ phút $- 7$ giờ $45$ phút $= 3$ giờ $45$ phút.\\
Chiều bắt đầu lúc: $11$ giờ $30$ phút $+ 1$ giờ $15$ phút $= 12$ giờ $45$ phút.\\
Chiều làm: $16$ giờ $45$ phút $- 12$ giờ $45$ phút $= 4$ giờ.\\
Tổng thời gian làm việc: $3$ giờ $45$ phút $+ 4$ giờ $= 7$ giờ $45$ phút (hoặc $16$ giờ $45$ phút $- 7$ giờ $45$ phút $- 1$ giờ $15$ phút $= 7$ giờ $45$ phút).\\
\textbf{Đáp án của câu hỏi là C}
}
\end{ex}

\textbf{Phần 2. Trắc nghiệm đúng sai}
\setcounterpageref{ex}{0}

\begin{ex} %[5D10H56]
Xét tính đúng sai của các phép đổi đơn vị đo thời gian sau:
\choiceTF
{\True $3$ năm $6$ tháng $= 42$ tháng.}
{$1{,}5$ ngày $= 30$ giờ.}
{\True $150$ giây $= 2{,}5$ phút.}
{\True $1$ giờ $15$ phút $= 1{,}25$ giờ.}
\loigiai{
\begin{itemize}
    \item $3 \times 12 + 6 = 42$ tháng. Đúng.
    \item $1{,}5 \times 24 = 36$ giờ $\neq 30$ giờ. Sai.
    \item $150 : 60 = 2{,}5$ phút. Đúng.
    \item $1 + 15 : 60 = 1{,}25$ giờ. Đúng.
\end{itemize}
\begin{itemize}
    \item a) \textbf{Đúng}.
    \item b) \textbf{Sai}.
    \item c) \textbf{Đúng}.
    \item d) \textbf{Đúng}.
\end{itemize}
}
\end{ex}

\begin{ex} %[5D10V56]
Xét tính đúng sai về lịch và thế kỉ:
\choiceTF
{\True Năm 2000 là năm nhuận và thuộc thế kỉ XX.}
{\True Năm 2026 không phải là năm nhuận và tháng 2 có 28 ngày.}
{Tháng 4, 6, 9, 11 là các tháng có $31$ ngày.}
{\True Quý 1 của một năm gồm các tháng 1, 2, 3.}
\loigiai{
\begin{itemize}
    \item Năm 2000 chia hết cho 400 nên là năm nhuận, và là năm cuối cùng của thế kỉ XX. Đúng.
    \item Năm 2026 không chia hết cho 4 nên không nhuận, tháng 2 có 28 ngày. Đúng.
    \item Tháng 4, 6, 9, 11 có 30 ngày (không phải 31 ngày). Sai.
    \item Quý 1 gồm tháng 1, tháng 2, tháng 3. Đúng.
\end{itemize}
\begin{itemize}
    \item a) \textbf{Đúng}.
    \item b) \textbf{Đúng}.
    \item c) \textbf{Sai}.
    \item d) \textbf{Đúng}.
\end{itemize}
}
\end{ex}

\textbf{Phần 3. Trả lời ngắn}
\setcounterpageref{ex}{0}

\begin{ex} %[5D10H56]
Đổi $3{,}2$ giờ ra phút.
\shortans{192}
\loigiai{
\textbf{Đáp số:} $192$\\
$3{,}2 \times 60 = 192$ phút.
}
\end{ex}

\begin{ex} %[5D10H56]
Đổi $4$ ngày $18$ giờ ra giờ.
\shortans{114}
\loigiai{
\textbf{Đáp số:} $114$\\
$4 \times 24 + 18 = 114$ giờ.
}
\end{ex}

\begin{ex} %[5D10V56]
Đổi $135$ phút ra giờ (dùng dấu phẩy nếu có).
\shortans{2,25}
\loigiai{
\textbf{Đáp số:} $2{,}25$\\
$135 : 60 = 2{,}25$ giờ.
}
\end{ex}

\begin{ex} %[5D10V56]
Đổi $5$ phút $24$ giây ra số thập phân có đơn vị là phút.
\shortans{5,4}
\loigiai{
\textbf{Đáp số:} $5{,}4$\\
$5 + 24 : 60 = 5{,}4$ phút.
}
\end{ex}

\begin{ex} %[5D10V56]
Một năm không nhuận có tất cả bao nhiêu giờ?
\shortans{8760}
\loigiai{
\textbf{Đáp số:} $8760$\\
$365 \times 24 = 8\,760$ giờ.
}
\end{ex}

\begin{ex} %[5D10C56]
Một người đi từ thành phố A lúc $6$ giờ $30$ phút và đến thành phố B lúc $9$ giờ $15$ phút. Trên đường người đó nghỉ giải lao $20$ phút. Tính thời gian thực tế người đó di chuyển trên đường theo phút.
\shortans{145}
\loigiai{
\textbf{Đáp số:} $145$\\
Thời gian từ lúc đi đến lúc đến: $9$ giờ $15$ phút $- 6$ giờ $30$ phút $= 2$ giờ $45$ phút $= 165$ phút.\\
Thời gian di chuyển: $165 - 20 = 145$ phút.
}
\end{ex}
